\chapter{Modélisation mathématique}
\label{chap:math}

Ce chapitre propose une formalisation ensembliste du système, de ses contraintes et de ses
indicateurs. Cette modélisation explicite les invariants que l'implémentation doit garantir
et fournit un langage rigoureux pour raisonner sur le matching et les performances.

\section{Ensembles fondamentaux}

Le système est décrit par les ensembles d'entités suivants, tous identifiés par des UUID :
\[
\begin{aligned}
&T \text{ (Tenants)}, \quad U \text{ (Users)}, \quad A \text{ (Actors)}, \quad
B \text{ (Businesses)}, \quad S \text{ (ServiceCatalogues)}, \\
&R \text{ (Resources)}, \quad Po \text{ (Portfolios)}, \quad SR \text{ (ServiceRequests)},
\quad I \text{ (Invoices)}, \quad BR \text{ (BusinessRules)}.
\end{aligned}
\]
L'ensemble des devises supportées est fini :
\[
C = \{\text{XAF, XOF, NGN, KES, GHS, USD, EUR, GBP}\}.
\]

\subsection{Décomposition des acteurs}
\[
A = A_{c} \cup A_{p}, \qquad A_{c} \cap A_{p} \neq \emptyset,
\]
où $A_c$ (consommateurs) et $A_p$ (fournisseurs) peuvent se recouper : un même acteur peut
à la fois consommer et fournir des services.

\subsection{Attributs des entités principales}
\begin{align*}
t  &= (\mathit{id}_T, \text{name}, \text{domain}, \text{isActive}, \text{createdAt}) \in T,\\
u  &= (\mathit{id}_U, \text{tenant}, \text{email}, \text{roles}, \text{isActive}, \dots) \in U,\\
a  &= (\mathit{id}_A, \text{user}, \text{role}, \text{bio}, \text{isActive}) \in A,\\
sr &= (\mathit{id}_{SR}, \text{tenant}, \text{consumer}, \text{provider}, \text{catalogue},
       \text{status}, \text{traceId}, \dots) \in SR,\\
i  &= (\mathit{id}_I, \text{serviceRequest}, \text{amount}, \text{currency}, \text{status},
       \text{issuedAt}, \text{paidAt}) \in I.
\end{align*}

\section{Variables de décision et fonction objectif}

\subsection{Variables}
Variable binaire de matching :
\[
x_{a_c, a_p, b} \in \{0,1\}, \quad \forall (a_c, a_p, b) \in A_c \times A_p \times B,
\]
valant $1$ si le consommateur $a_c$ est mis en relation avec le fournisseur $a_p$ pour le
métier $b$. Variable d'état d'une demande :
\[
\sigma_{sr} \in \Sigma_{SR} = \{\text{PENDING, ACCEPTED, IN\_PROGRESS, FULFILLED, CANCELLED}\}.
\]
\iffalse
\subsection{Objectif multi-critère}
On cherche à maximiser le nombre de transactions réussies et le chiffre d'affaires, tout en
minimisant la durée de traitement :
\[
\max\; \alpha \sum_{sr \in SR} \mathbb{1}_{\sigma_{sr} = \text{FULFILLED}}
     + \beta  \sum_{i \in I} \text{amount}(i)\,\mathbb{1}_{\sigma_i = \text{PAID}}
     - \gamma \sum_{sr \in SR} \text{duration}(sr),
\]
\[
\text{avec } \alpha + \beta + \gamma = 1, \quad \alpha,\beta,\gamma \geq 0.
\]
\fi

\section{Contraintes du système}

\subsection{Contraintes de cardinalité}
\begin{align}
&\forall u_1, u_2 \in U,\ u_1 \neq u_2 : \ u_1.\text{email} \neq u_2.\text{email}
   &&\text{(unicité de l'email)}\\
&\forall a \in A,\ \exists!\, po \in Po : po.\text{actor} = a
   &&\text{(Actor} \to \text{Portfolio, 1:1)}\\
&\forall po \in Po : \ po.\text{businesses} \subseteq B
   &&\text{(Portfolio} \to \text{Business, N:M)}\\
&\forall sr \in SR,\ \exists!\, i \in I : i.\text{serviceRequest} = sr
   &&\text{(ServiceRequest} \to \text{Invoice, 1:1)}
\end{align}

\subsection{Contraintes métier}
Validité d'une demande :
\[
\text{valid}(sr) \iff
\begin{cases}
sr.\text{consumer}.\text{role} = \text{CONSUMER},\\
sr.\text{provider}.\text{role} = \text{PROVIDER},\\
sr.\text{provider}.\text{isActive} = \text{vrai}.
\end{cases}
\]
Compatibilité de matching :
\[
\text{compatible}(a_c, a_p, b) \iff \exists\, po \in Po :
\big(po.\text{actor} = a_p \ \wedge\ b \in po.\text{businesses}\big).
\]
Contrainte de devise :
\[
\forall i \in I : i.\text{currency} \in C, \qquad \forall s \in S : s.\text{currency} \in C.
\]
Contrainte d'isolation multi-tenant (invariant de sécurité) :
\[
\forall sr \in SR : sr.\text{tenant} = \tau \implies
\big(sr.\text{consumer}.\text{user}.\text{tenant} = \tau \ \wedge\
      sr.\text{catalogue}.\text{business}.\text{tenant} = \tau\big).
\]

\subsection{Contrainte de workflow : machine à états}
Soit $\delta_{SR} : \Sigma_{SR} \times \mathcal{A} \to \Sigma_{SR}$ la fonction de
transition, avec $\mathcal{A} = \{\text{create, accept, start, fulfill, cancel}\}$. Les
transitions valides sont données au tableau~\ref{tab:fsm-math}.

\begin{table}[H]
\centering
\caption{Transitions valides de la machine à états \texttt{ServiceRequest}}
\label{tab:fsm-math}
\renewcommand{\arraystretch}{1.2}\footnotesize
\begin{tabular}{lll}
\toprule
\textbf{État courant} & \textbf{Action} & \textbf{État suivant} \\
\midrule
$\emptyset$    & create  & PENDING \\
PENDING        & accept  & ACCEPTED \\
PENDING        & cancel  & CANCELLED \\
ACCEPTED       & start   & IN\_PROGRESS \\
ACCEPTED       & cancel  & CANCELLED \\
IN\_PROGRESS   & fulfill & FULFILLED \\
IN\_PROGRESS   & cancel  & CANCELLED \\
\bottomrule
\end{tabular}
\end{table}

\section{Algorithmes et complexité}

\subsection{Recherche des fournisseurs d'un métier}
\begin{lstlisting}[language={},caption={FindProvidersByBusiness},captionpos=b]
Entree : b (business cible)
Sortie : ensemble des providers compatibles
1: Resultat <- vide
2: pour chaque po dans Po :
3:     si b appartient a po.businesses :
4:         a <- po.actor
5:         si a.role = PROVIDER et a.isActive : Resultat <- Resultat U {a}
6: retourner Resultat
\end{lstlisting}
\textbf{Complexité :} $O(|Po|)$ dans le pire cas, $O(1)$ amorti avec un index sur les
métiers du portfolio.

\subsection{Recherche par filtres combinés}
Pour un ensemble de filtres $F = \{f_1, \dots, f_n\}$, le résultat est l'intersection
\[
\text{result}(F) = \bigcap_{f \in F} \text{applyFilter}(f, \text{Univers}),
\]
optimisée en appliquant d'abord le filtre le plus sélectif :
\[
\text{selectivity}(f) = \frac{|\text{applyFilter}(f, \text{Univers})|}{|\text{Univers}|},
\qquad \text{ordre croissant de } \text{selectivity}.
\]

\section{Indicateurs et métriques (KPIs)}

Sur une période $[t_1, t_2]$, on définit :
\[
\boxed{T_{\text{conv}} = \frac{N_{\text{fulfilled}}}{N_{\text{requests}}} \times 100\%},
\qquad
\boxed{PA = \frac{CA}{N_{\text{paid}}}},
\]
où le chiffre d'affaires est
\[
CA(t_1, t_2) = \sum_{\substack{i \in I \,:\, \sigma_i = \text{PAID}\\ i.\text{paidAt} \in [t_1,t_2]}}
i.\text{amount},
\]
et le temps moyen de réalisation
\[
\boxed{T_{\text{compl}} = \frac{1}{N_{\text{fulfilled}}}
\sum_{\substack{sr \in SR \,:\, \sigma_{sr} = \text{FULFILLED}}}
\big(sr.\text{fulfilledAt} - sr.\text{requestedAt}\big)}.
\]
La part d'un métier $b$ dans le chiffre d'affaires s'écrit
$\text{Share}_b = CA_b / CA \times 100\%$, et la croissance mensuelle
$\text{Growth}_{\text{MoM}}(t) = (CA(t) - CA(t-1)) / CA(t-1) \times 100\%$. Ces indicateurs
sont calculés par le service d'analytique 
% ════════════════════════════════════════════════════════
% À COLLER DIRECTEMENT APRÈS LA DERNIÈRE LIGNE DU 5.5
% Aucun package ou couleur supplémentaire requis.
% ════════════════════════════════════════════════════════

\section{Architecture en couches : composition de fonctions}

Ce modèle formalise le traitement d'une requête à travers les cinq
couches de BCaaS. Chaque couche est représentée par une fonction de
transformation, et le traitement complet est leur composée.

\subsection{Définition formelle}

On associe à chaque couche $i \in \{1,\ldots,5\}$ une fonction :
\begin{align*}
f_1 &: r   \mapsto r_1 && \text{Infrastructure (connexion BD, cache)}\\
f_2 &: r_1 \mapsto r_2 && \text{Transport \& Messaging (REST ou Kafka)}\\
f_3 &: r_2 \mapsto r_3 && \text{Tenant \& Routing (résolution \texttt{tenant\_id})}\\
f_4 &: r_3 \mapsto r_4 && \text{Context \& Policy (JWT, RBAC, SLA)}\\
f_5 &: r_4 \mapsto r_5 && \text{Business Capabilities (logique métier)}
\end{align*}

Le traitement complet d'une requête $r$ est la \textbf{composée} de
ces cinq fonctions :
\[
\text{BCaaS}(r)
\;=\;
(f_5 \circ f_4 \circ f_3 \circ f_2 \circ f_1)(r)
\;=\;
f_5\!\Bigl(f_4\!\bigl(f_3\!\bigl(f_2\!\bigl(f_1(r)\bigr)\bigr)\bigr)\Bigr).
\]

\subsection{Propriété : encapsulation stricte}

Pour tout $i \in \{1,\ldots,5\}$, la fonction $f_i$ ne reçoit en entrée
que $r_{i-1}$ et est indépendante de $f_j$ pour tout $j > i$ :
\[
f_i \;:\; \mathcal{R}_{i-1} \;\longrightarrow\; \mathcal{R}_i,
\qquad f_i \perp f_j \text{ pour } j > i.
\]
Modifier $f_5$ n'affecte jamais $f_1$ : c'est la traduction formelle du
principe d'encapsulation du modèle OSI.

\subsection{Exemple}

Considérons une \texttt{ServiceRequest} de création de métier
(\texttt{POST /api/businesses}) émise par un acteur authentifié.
La requête $r$ traverse les cinq couches :
$f_1$ ouvre la connexion PostgreSQL,
$f_2$ sélectionne le transport REST synchrone,
$f_3$ résout le \texttt{tenant\_id} depuis l'en-tête \texttt{X-Tenant-ID},
$f_4$ valide le token JWT et le rôle \texttt{ADMIN},
et $f_5$ persiste l'entité et retourne \texttt{201 Created} :
\[
\text{BCaaS}(r) = f_5\bigl(f_4\bigl(f_3\bigl(f_2\bigl(f_1(r)\bigr)\bigr)\bigr)\bigr)
= r_5 = \texttt{201 Created}.
\]
Si $f_4$ échoue (JWT expiré), la composition s'arrête en $r_4$ et
retourne \texttt{401 Unauthorized}, sans jamais atteindre $f_5$.


\section{Isolation multi-tenant : théorie des ensembles}

Ce modèle formalise l'étanchéité des données entre les tenants de
l'ensemble $T$ (défini à la section~5.1). Il constitue la preuve
ensembliste globale de la contrainte d'isolation posée à la section~5.3.

\subsection{Définitions}

On note $\mathcal{E}$ l'ensemble de toutes les lignes présentes en base
de données (tables $U$, $A$, $B$, $SR$, $I$, etc. confondues).
On définit la fonction d'appartenance $\pi : \mathcal{E} \to T$ par :
\[
\pi(e) = \tau_k
\quad \Longleftrightarrow \quad
\text{la colonne \texttt{tenant\_id} de } e \text{ vaut } \tau_k.
\]
L'ensemble des données du tenant $\tau_k$ est :
\[
\mathcal{D}(\tau_k) \;=\; \bigl\{\, e \in \mathcal{E} \;\big|\; \pi(e) = \tau_k \,\bigr\}.
\]

\subsection{Propriété : isolation totale}

\[
\forall\, \tau_i, \tau_j \in T,\quad
\tau_i \neq \tau_j
\;\implies\;
\mathcal{D}(\tau_i) \cap \mathcal{D}(\tau_j) = \emptyset.
\]

Les ensembles $\mathcal{D}(\tau_1), \ldots, \mathcal{D}(\tau_n)$
forment une \textbf{partition} de $\mathcal{E}$ :
\[
\mathcal{E}
\;=\;
\biguplus_{k=1}^{n} \mathcal{D}(\tau_k)
\;=\;
\mathcal{D}(\tau_1) \;\uplus\; \mathcal{D}(\tau_2)
\;\uplus\; \cdots \;\uplus\; \mathcal{D}(\tau_n).
\]
Aucune entité n'est « sans tenant » ni « dans deux tenants à la fois ».
Cette propriété est garantie dans le code par le \texttt{TenantFilter}
de Spring Security, qui injecte le \texttt{tenant\_id} dans chaque
requête SQL via un \texttt{ThreadLocal}.

\subsection{Exemple}

Deux tenants sont inscrits : $\tau_1$ (Auto-École Centrale) et
$\tau_2$ (Clinique Mélen). Les \texttt{ServiceRequests} associées sont :
\[
\mathcal{D}(\tau_1) = \{sr_1, sr_2, sr_3\},
\qquad
\mathcal{D}(\tau_2) = \{sr_4, sr_5\},
\qquad
\mathcal{D}(\tau_1) \cap \mathcal{D}(\tau_2) = \emptyset.
\]
Un acteur de $\tau_2$ exécutant \texttt{GET /api/service-requests} ne
verra jamais $sr_1$, $sr_2$ ou $sr_3$ : le \texttt{TenantFilter} ajoute
automatiquement \texttt{WHERE tenant\_id = '$\tau_2$'} à chaque requête.


\section{Provisionnement d'un tenant : suite finie d'états}

Ce modèle formalise l'inscription d'un nouveau tenant comme une suite
d'états atomiques. Il est distinct de la machine à états des
\texttt{ServiceRequests} (section~5.3.3) : il ne s'exécute qu'une seule
fois par tenant, au moment de son inscription sur la plateforme.

\subsection{Les sept états du provisionnement}

Soit $\mathcal{S} = (S_0, S_1, \ldots, S_6)$ la suite des états :

\begin{table}[H]
\centering
\caption{Suite d'états du provisionnement d'un tenant}
\label{tab:prov-states}
\renewcommand{\arraystretch}{1.2}\small
\begin{tabular}{lll}
\toprule
\textbf{État} & \textbf{Nom} & \textbf{Description et précondition} \\
\midrule
$S_0$ & Requête reçue
  & Réception de \texttt{POST /api/tenants \{name, domain, sla\}}. \\
$S_1$ & Validation du domaine
  & \textbf{Précondition :}
    $\text{dom}_{\text{req}} \notin \{t.\text{domain} \mid t \in T\}$.
    Sinon : \texttt{409 Conflict}. \\
$S_2$ & Génération de l'id
  & Génération d'un $\text{id} \in \text{UUID}_4$
    (collision $\approx 10^{-36}$). \\
$S_3$ & Insertion en base
  & \texttt{INSERT INTO tenant} en transaction atomique.
    Rollback si échec. \\
$S_4$ & Contexte d'isolation
  & Initialisation $\mathcal{D}(t_{\text{new}}) \leftarrow \emptyset$,
    \texttt{TenantFilter} configuré. \\
$S_5$ & Génération JWT
  & Token signé $(\text{tenant\_id}, \text{roles}, \exp)$,
    $\exp = t_{\text{now}} + 24\,\text{h}$. \\
$S_6$ & Tenant actif
  & $t_{\text{new}} \in T$ et
    $\mathcal{D}(t_{\text{new}}) = \emptyset$. \\
\bottomrule
\end{tabular}
\end{table}

\subsection{Propriété : injectivité du domaine}

La précondition de $S_1$ garantit que la fonction
$\text{domain} : T \to \mathit{Dom}$ est \textbf{injective} :
\[
\text{domain}(t_i) = \text{domain}(t_j) \;\implies\; t_i = t_j.
\]
Deux tenants distincts ne peuvent jamais partager le même domaine.
Cette propriété complète les contraintes de cardinalité de la
section~5.3 et est analogue à l'unicité des adresses IP sur un
réseau : elle permet un routage sans ambiguïté par $f_3$ (section~5.6).

\subsection{Exemple}

Un gestionnaire soumet
\texttt{POST /api/tenants \{"domain":"cabinet-dupont"\}}.
En $S_1$, le système vérifie
$\texttt{"cabinet-dupont"} \notin \{t.\text{domain} \mid t \in T\}$ :
précondition satisfaite.
En $S_2$, l'UUID \texttt{b9c2-4f1a-\ldots} est généré.
En $S_3$, la ligne est insérée de façon atomique.
En $S_4$, $\mathcal{D}(t_{\text{new}}) \leftarrow \emptyset$.
En $S_5$, le JWT initial est signé.
En $S_6$ :
\[
t_{\text{new}} \in T, \quad
\mathcal{D}(t_{\text{new}}) = \emptyset, \quad
\text{domain}(t_{\text{new}}) = \texttt{"cabinet-dupont"}.
\]
Si une seconde requête tentait le même domaine, $S_1$ retournerait
\texttt{409 Conflict}, préservant l'injectivité.
(chapitre~\ref{chap:observ}).
